Experiment~2 indicated that Proposed-Med, evaluated separately at Fp1 and at Fp2, the two electrodes identified in Experiment~1 as consistently the best-performing, was insensitive to epoch duration in the pooled analysis: macro-$F_1$ spanned only 0.41 (Fp1) and 0.50 (Fp2) percentage points across the seven tested durations, and none differed significantly from the 30-second reference after correction for either electrode (Figure~\ref{fig:f1_by_epoch}). Cao2018 was similarly stable for both electrodes, with $F_1$ spreads of 0.38 (Fp1) and 0.33 (Fp2) percentage points and no significant difference at any duration. This pattern is useful because epoch duration is commonly determined by the experimental paradigm rather than by the blink detector. The corpus-specific results nevertheless qualify the pooled finding, and reveal an electrode-dependent asymmetry: on Internal, Fp1's $F_1$ declined from 86.32\% at 30 seconds to 85.92\% at 50 seconds ($p=0.002$), 85.67\% at 60 seconds ($p<0.001$), and 85.04\% at 120 seconds ($p<0.001$), all three remaining significant after correction, whereas Fp2 declined significantly only at 120 seconds (86.82\% to 85.57\%, $p<0.001$), with its 50- and 60-second drops falling short of significance ($p=0.080$ and $p=0.064$). Thus, the evidence supports stable pooled performance and no detected change from 10 to 40 seconds in either corpus for either electrode, but it does not establish equivalence across the complete 10--120-second range within every dataset, and the point at which Internal performance becomes unstable is itself electrode-dependent. Longer epochs may alter the composition of the signal used during screening and threshold estimation, but the present analysis did not test that mechanism directly.
